\chapter{不确定度、可辨识性与模型检验}
\section{信息矩阵与条件协方差}
先考虑无化学限制、局部线性且误差模型成立的情况。独立变量的衍射 Fisher 信息近似为
\begin{equation}
\mat H_D=\mat J_D^{\mathsf T}\mat J_D,
\qquad \Cov(\widehat{\vct z})\simeq\mat H_D^{-1}.
\end{equation}
这里 $\mat J_D$ 已含观测标准差。加入限制后，联合曲率为 $\mat H=\mat H_D+\mat J_R^{\mathsf T}\mat J_R$。若将限制解释为具有指定误差的独立辅助观测，$\mat H^{-1}$ 对应联合观测的局部协方差；若将限制解释为高斯先验，它是局部后验曲率近似。对固定限制目标、只重复采集衍射数据的频率学抽样，线性化协方差则为
\begin{equation}
\mat C_{\mathrm{repeat}}\simeq\mat H^{-1}\mat H_D\mat H^{-1}.
\end{equation}
这几种解释不能混用。源码返回的是联合曲率的逆，并以给定限制、误差和模型为条件，不声称自动完成了频率学重复采样或完整 Bayesian 推断。

如果输入标准差仅确定到共同倍数，可选择再以残差方差估计缩放矩阵；如果标准差本来就是绝对量，额外乘以拟合优度可能缩小或扩大不确定度。软件把两种数值选择记录为 \code{absolute} 与 \code{reduced_chi2}，不在结果生成时暗中切换。包含限制时，还应明确这种整体缩放是否符合限制的误差解释。

对物理参数 $\vct p=\mathcal P(\vct z)$，一阶误差传播为
\begin{equation}
\mat C_p=\mat D\mat C_z\mat D^{\mathsf T},\qquad
\mat D=\frac{\partial\mathcal P}{\partial\vct z}.
\end{equation}
因此，logit、Cholesky 元素和矩阵对数的标准差不应直接标在占有率或 $U_{ij}$ 上。原生结果分别保存优化变量协方差和传播后的物理量标准不确定度。

\section{SVD 与秩缺失}
对 Jacobian 作奇异值分解
\begin{equation}
\mat J=\mat U\mat S\mat V^{\mathsf T},\qquad
\mat S=\diag(s_1,\ldots,s_r,0,\ldots).
\end{equation}
当存在 $\mat J\vct v=0$ 的非零方向，数据在一阶上不能区分沿该方向的变化。对可辨识子空间可计算
\begin{equation}
\mat H^+=\mat V_r\diag(s_i^{-2})\mat V_r^{\mathsf T},
\end{equation}
但它不是所有参数的有限协方差。伪逆把零信息方向的逆值设为零；若直接取其对角线开方，会把“完全不知道”误写成“误差为零”。实现出现秩缺失时返回空的完整协方差，另行报告伪逆、奇异值与零空间。

以下三种零方差必须区分：固定参数没有参与估计；原点规范移除了不可观测的公共方向；观测真正给出了高精度估计。前两者的标准不确定度在机器结果中标为未估计，不用数值零冒充测量精度。

\section{数据与限制的识别作用}
令 $\mat J_D$ 为衍射残差 Jacobian，$\mat J_R$ 为限制 Jacobian，则
\begin{equation}
\mat H=\mat J_D^{\mathsf T}\mat J_D+\mat J_R^{\mathsf T}\mat J_R.
\end{equation}
限制可以补足 $\mat J_D$ 的缺失方向，使联合矩阵满秩。这是合理的统计建模，但结果必须说明某些组合主要由限制决定。软件同时返回衍射 Jacobian 秩和联合 Jacobian 秩，便于 agent 识别这种情况。

限制过强还可能把错误的化学模型拟合成几何整齐的结构。检验应同时考察限制残差、衍射残差和去除该限制后参数变化。一次优化的低目标值不能代替竞争模型比较。

\section{R 因子与拟合优度}
令 $F_o=\sqrt{\max(y,0)}$，$F_c^{\mathrm{scaled}}=\sqrt{\mu}$。常用指标为
\begin{align}
R_1&=\frac{\sum\left|F_o-F_c^{\mathrm{scaled}}\right|}{\sum F_o},\\
wR_2&=\sqrt{\frac{\sum_hw_h(y_h-\mu_h)^2}{\sum_hw_hy_h^2}},\\
S&=\sqrt{\frac{\sum_hw_h(y_h-\mu_h)^2}{n_{\mathrm{work}}-n_{\mathrm{par}}}}.
\end{align}
原生结果分别列出全部工作反射的 $R_1$ 与 $y>2\sigma$ 子集的 $R_1$。子集只用于报告，不从精修目标中剔除弱反射。拟合优度的分母明确采用衍射观测数与独立参数数，限制数和限制平方和单独报告。存在强限制、活动边界或秩缺失时，这个量不再具有通常无约束线性模型的 $\chi^2$ 标定，应作为描述性诊断解释。

同一个 $R_1$ 值可能对应不同分辨率、数据完整度和反射筛选规则，所以跨数据集比较必须同时提供这些条件。低 $R$ 不证明原子类型、绝对结构或无序模型唯一正确。

\section{留出反射与信息泄漏}
随机逐行分割反射可能使同一对称等价组的一部分进入训练集，另一部分进入测试集。若数据包含重复观测、Friedel 对或孪晶关联，这种拆分会产生泄漏。软件以对称和孪晶生成的反射轨道为单位分组，并把同组全部分配到同一集合。

固定工作集与测试集后，参数估计只使用工作集。测试指标用于评价外推到未用于估计的观测。若 agent 反复查看同一测试集并据此调整空间群、无序模型或限制强度，该集合就已经参与模型选择；最终的独立评估需要另一个预留集合或外部数据。测试集不是可无限重复使用的真实性证明。

有限数据的小分子结构未必适合留出很大比例的反射。方案应根据独立反射数与假设复杂度确定，并保留全数据最终精修与独立检验的区别。公开 FeN$_4$ 例子只有 71 条反射，本报告使用全数据复现归档精修，不另外作不足以稳定评价的随机留出。

\section{模型复杂度与残余密度}
增加原子、放开占有率或引入额外孪晶分数，通常能降低工作残差。只有当下降与新增参数的可辨识性、残余密度消除及独立数据一致时，复杂模型才获得额外支持。局部信息准则可以提供线索，但在不同权重、限制或边界条件下，其简单参数计数并不总成立。

差值密度中的峰可能来自缺失原子，也可能来自错误位移参数、吸收误差、孪晶或截断波纹。一个合理的 agent 应记录峰到原子的位置关系、正负峰成对性、分辨率依赖和替代解释，再提出可检验的模型改变。每次改变都形成可回溯的分支，而不是覆盖唯一的结构文件。
